\subsection{Nonlinear redundancy with no inactive retained coordinate}
\label{sec:nonlinearhub}
The distinction is not limited to redundant linear combinations.
Consider the frozen four-qubit, one-layer Hubregtsen source circuit
\cite{hubregtsen2022,squlearn}
with two input features and local Z readout. It prepares $|+\rangle$
on each qubit, applies $R_z(x_{q\bmod2})$ and $R_y(\theta_q)$,
and ends with controlled-Z rotations carrying four further symbols.
Those controlled rotations commute with every measured Z and are
already removed by the strengthened observable-commutant baseline.
The four retained features are exactly
\begin{equation}
 f_q(x)=-\sin\theta_q\cos x_{q\bmod2},\qquad q=0,\ldots,3.
 \label{eq:hubfeatures}
\end{equation}
Exact Heisenberg expansions of the frozen compiled IR reproduce
each identity, including its sign and source-parameter index.

\begin{proposition}[A nonlinear two-coordinate factorization]
For this source configuration and unit-weight local Z readout,
all four retained source coordinates are individually kernel-active.
The minimum dimension of a smooth kernel factorization and the
generic differential rank are both two, whereas any factorization
through a constant linear map of the retained coordinates requires
rank four.
\end{proposition}
\begin{proof}
Let $c_j=\cos x_j-\cos x'_j$. Equation~\eqref{eq:hubfeatures} gives
\begin{align}
 d^2_\theta(x,x')&=\beta_0 c_0^2+\beta_1 c_1^2,\label{eq:hubdistance}\\
 \beta_0&=\sin^2\theta_0+\sin^2\theta_2,\nonumber\\
 \beta_1&=\sin^2\theta_1+\sin^2\theta_3.\label{eq:hubbeta}
\end{align}
This is a global smooth factorization through two real variables.
Set the four retained angles to $\pi/4$ and choose input pairs
$((0,0),(\pi/4,0))$ and $((0,0),(0,\pi/4))$ in $(-1,1)^2$.
The distance Jacobian minor in source columns $\theta_0,\theta_1$
is diagonal with entries $(1-1/\sqrt2)^2$, hence
\begin{equation}
 \det J=(1-1/\sqrt2)^4=\frac{17}{4}-3\sqrt2>0.
 \label{eq:hubminor}
\end{equation}
Analyticity and the chain rule therefore give generic rank and
minimum smooth factor dimension two. Each source derivative
$\sin(2\theta_q)c_{q\bmod2}^2$ is not identically zero, so every
retained coordinate is active. A constant null vector $v$ would
require $v_0\sin2\theta_0+v_2\sin2\theta_2=0$ and
$v_1\sin2\theta_1+v_3\sin2\theta_3=0$ for all independently
variable angles. Each coefficient must vanish. A rank-deficient
constant linear factor would supply a nonzero such vector and is
therefore impossible.
\end{proof}

Here an exact constant-null-space calculation returns the valid
upper bound four, not the minimal nonlinear dimension two. This
source-derived example delimits the completeness of that upper-bound
method; it does not establish a general nonlinear reduction algorithm.
The two input features are reused across four qubits, and the final
controlled rotations are invisible to the selected readout. The
factorization is thus a readout-specific reduction to a weighted
classical cosine-feature kernel, not evidence of quantum advantage
or a newly efficient hardware compiler. Nor does ordinary gradient
descent in $\beta$ reproduce the original finite-step optimizer in
$\theta$ without a separately specified update map.
